\noindent\textit{Enunciado adaptado de los problemas 4 y 5 del Banco de Problemas de Aprendizaje por Refuerzo (Modelos de Decisión en el Tiempo, Módulo 3, IIND-3518, Universidad de los Andes).}

\bigskip

\noindent Una tienda mantiene un inventario $S_t \in \{0, 1, 2\}$. Al comenzar el día elige un pedido $A_t \in \{0, \dots, 2 - S_t\}$, que llega de inmediato. La demanda del día es $D_t \in \{0, 1\}$, con probabilidad $1/2$ para cada valor. Las ventas, el inventario siguiente y la recompensa son
\[
Y_t = \min\{S_t + A_t,\ D_t\}, \qquad S_{t+1} = S_t + A_t - Y_t, \qquad R_{t+1} = 5Y_t - 2A_t - S_{t+1}.
\]
Es decir, cada unidad vendida deja 5, cada unidad pedida cuesta 2 y cada unidad que queda en inventario al final del día cuesta 1. El negocio opera indefinidamente, empieza con $S_0 = 1$ y usa $\gamma = 0{,}9$.

\begin{enumerate}
\item Formule el problema para Q-learning: estados, acciones factibles y qué observa el algoritmo después de cada decisión.
\item Entrene Q-learning con $\alpha = 0{,}005$, $\varepsilon = 0{,}2$, $10^6$ pasos y semilla 2026. Presente la tabla $Q$ y la política greedy aprendida.
\item Calcule la política óptima de forma exacta con \emph{value iteration} y compárela con la aprendida en el inciso anterior. Compare también la ganancia máxima de cada estado con la ganancia de seguir la política aprendida.
\item Interprete la política óptima en cada estado.
\end{enumerate}

\noindent\textbf{Solución:} \texttt{qlearning\_problem\_solutions/qlearning\_tienda.ipynb}.
